% Рисунок из тетради (лемма Дюбуа--Реймона): K ⊂ Ω̃ ⊂⊂ Ω.
\begin{figure}[htbp]
\centering
\begin{tikzpicture}[scale=1.0]
  % внешняя область Omega (может быть неограниченной)
  \draw[thick, fill=blue!3]
    plot[smooth cycle, tension=0.8]
      coordinates {(0,0) (2.6,1.0) (5.0,0.3) (5.6,-1.6) (3.4,-3.0) (0.8,-2.6) (-0.9,-1.2)};
  \node at (4.8,-0.4) {$\Omega$};

  % промежуточная ограниченная область Omega tilde
  \draw[thick, dashed, fill=blue!8]
    plot[smooth cycle, tension=0.85]
      coordinates {(1.0,-0.4) (2.6,0.2) (3.9,-0.6) (3.6,-2.0) (2.0,-2.4) (0.9,-1.5)};
  \node at (3.35,-1.95) {$\widetilde{\Omega}$};

  % компакт K
  \filldraw[pattern=north east lines, pattern color=blue!45, draw=blue!60!black, thick]
    plot[smooth cycle, tension=0.9]
      coordinates {(1.6,-0.9) (2.6,-0.75) (3.0,-1.4) (2.3,-1.9) (1.5,-1.6)};
  \node at (2.25,-1.35) {$K$};
\end{tikzpicture}
\caption{К доказательству леммы: по компакту $K \subset \Omega$ выбирается ограниченная
открытая $\widetilde{\Omega}$ с $K \subset \widetilde{\Omega} \Subset \Omega$. На ней
$f \in L^1(\widetilde{\Omega})$, и там же живут усреднения $g_\rho$.}
\label{fig:nested}
\end{figure}
